\chapter{Security, Cost and Deployment}
\label{ch:security}

\begin{quote}\itshape
Content a tool fetched is data. It is never an instruction, no matter how
confidently it is phrased.
\end{quote}

\noindent
An agent is a system that reads untrusted text and then decides what to do next.
Stated that way the security problem is obvious and the mitigations follow
directly --- and the single organising principle is that \textbf{every effective
control is enforced outside the model}.

% ============================================================
\section{Prompt injection}
\label{sec:injection}
\index{security!prompt injection}

\mermaidfig{sec-injection-path}{%
  The path, and four controls. None of them asks the model to decline.}{%
  Prompt injection and the controls that stop it.}

A web page your agent reads, a ticket comment, an email, a PDF --- any of them
can contain text addressed to the model rather than to the user. The model has no
reliable way to distinguish the two, because to a model they are the same
sequence of tokens.

\begin{warning}
Do not attempt to solve this by instructing the model to ignore instructions in
tool output. It helps a little, it fails under adversarial pressure, and relying
on it means your control is a preference. Write the instruction if you like ---
and put a real control behind it.
\end{warning}

\begin{center}
\small
\begin{tabularx}{\linewidth}{@{}lX@{}}
\toprule
\textbf{Control} & \textbf{What it actually does} \\
\midrule
Tool allowlist per path & The retrieval path physically cannot reach
  \code{send\_email} \\[3pt]
Authorisation in the wrapper & Checked in code against the \emph{user's}
  permissions (\S\ref{sec:authz}) \\[3pt]
Approval gate & A human sees the recipient before it goes
  (Ch.~\ref{ch:interrupts}) \\[3pt]
Egress policy & The sandbox cannot reach the internet, so exfiltration has
  nowhere to go \\
\bottomrule
\end{tabularx}
\end{center}

% ============================================================
\section{Authorisation}
\label{sec:authz}
\index{security!tool authorisation}

The rule: \textbf{a tool runs with the user's permissions, not the service's.}

\begin{python}[caption={The check, in code the model cannot reach}, label={lst:authz-mw}]
class ToolAuthorisation(AgentMiddleware):
    def wrap_tool_call(self, request, handler):
        if not permitted(request.runtime.context.user_id,
                         request.tool_call["name"]):
            return ToolMessage(
                content="Not permitted.",
                tool_call_id=request.tool_call["id"],
            )                                # handler never called
        return handler(request)
\end{python}

Chapter~\ref{ch:agents} showed the mechanism; this is why it matters. Returning
without calling \code{handler} short-circuits the pipeline, and the refusal comes
back as a tool message so the model can explain it rather than the run failing.

The identity comes from \code{runtime.context} --- Chapter~\ref{ch:stategraph}
insisted it belongs there rather than in state, and this is the reason: state is
mutable by any node and is checkpointed, so a user id in state is a permission
that can be overwritten mid-run and then persisted.

% ============================================================
\section{Personal data}
\label{sec:pii}
\index{security!PII}

Three obligations, and the third is the one people miss.

\textbf{Redaction before the provider sees it.} Chapter~\ref{ch:agents}'s PII
middleware, with the default corrected: it is input-only unless you ask for
\code{apply\_to\_output} and \code{apply\_to\_tool\_results}. Data flowing out of
a tool into the context is not covered by the default.

\textbf{Processing region.} Which country the inference happens in is a
contractual question with a technical answer, and providers differ in what they
will commit to.

\textbf{Checkpoints are personal data.} This is the one that gets missed.
Chapter~\ref{ch:persistence}'s checkpointer writes the entire conversation to
your database, every superstep. It is a complete record of what a user said, and
it needs a retention policy, a deletion path for subject-access requests, and to
be in scope for whatever review your other personal data gets.

\begin{warning}
\textbf{A checkpoint is not an audit log.} They look similar and answer different
questions. A checkpoint is state for resumption --- it is overwritten, forked by
time travel (Ch.~\ref{ch:interrupts}), and deleted by retention. An audit log is
append-only, immutable, and retained deliberately.

If you need to answer ``who did what, with which tool, and what came back'',
write that separately in \code{wrap\_tool\_call}. Deriving it from checkpoints
will fail the first time someone asks a question about a deleted thread.
\end{warning}

% ============================================================
\section{Cost}
\label{sec:cost-control}
\index{cost!budget}

Chapter~\ref{ch:context} measured 5$\times$ between context strategies and
Chapter~\ref{ch:multi-agent} measured 2.8$\times$ between topologies. Those are
the two large levers, and they are design decisions rather than settings. What
remains is enforcement:

\begin{center}
\small
\begin{tabularx}{\linewidth}{@{}lX@{}}
\toprule
\textbf{Control} & \textbf{Where} \\
\midrule
Model per role & Cheap for routing, strong for work
  (Ch.~\ref{ch:multi-agent}) \\[3pt]
Call limits & \code{thread\_limit} and \code{run\_limit}
  (Ch.~\ref{ch:agents}) \\[3pt]
Prompt caching & Stable prefix first (Ch.~\ref{ch:context}) \\[3pt]
Node caching & \api{CachePolicy} plus a cache at compile time
  (Ch.~\ref{ch:persistence}) \\[3pt]
Budget cap & A running total in state; halt when exceeded \\
\bottomrule
\end{tabularx}
\end{center}

\begin{note}
Alert on the budget, not on the invoice. A per-conversation cost recorded from
\code{usage\_metadata} (Ch.~\ref{ch:models}) turns a monthly surprise into a
threshold you can page on --- and Chapter~\ref{ch:observability}'s CI check turns
it into something a pull request can fail.
\end{note}

% ============================================================
\section{Sandboxing}
\label{sec:sandbox}
\index{security!sandbox}

If your agent executes code --- and the pressure to add that is considerable ---
the model is now writing programs on your infrastructure from untrusted input.

Three layers, all of them necessary: a container with no network, a read-only
filesystem, dropped capabilities and a non-root user; resource limits, because an
infinite loop is a denial of service; and an egress policy, because without
network access an injection has nowhere to send what it stole.

\begin{warning}
Syntax-tree validation is not a security boundary. Neither is a denylist of
dangerous functions. Both are worth having and neither is what stands between you
and a compromised host --- that is the container and the network policy.
\end{warning}

% ============================================================
\section{Deployment}
\label{sec:deployment}

\mermaidfig{deploy-shutdown}{%
  SIGTERM to exit, and the one thing that breaks all of it.}{%
  Graceful shutdown sequence under Kubernetes.}

Chapter~\ref{ch:serving} built the service; this is what surrounds it.

\textbf{The container} runs as non-root, on a slim base, with dependencies from
the lockfile rather than resolved at build time --- so the image is what you
tested.

\textbf{Probes.} Readiness should fail the moment shutdown begins, so the load
balancer stops sending work while the pod is still serving what it has. Liveness
should not depend on the provider being up, or a provider outage restarts every
pod you own at the moment they are least able to recover.

\textbf{The grace period must exceed your longest run}, or the drain in the
diagram is decorative. If your runs are longer than a sensible grace period, that
is Chapter~\ref{ch:serving}'s queue telling you where they belong.

\begin{warning}
Every part of the shutdown sequence depends on cancellation propagating. One
swallowed \code{CancelledError} anywhere in the request path
(Ch.~\ref{ch:cancellation}) turns a graceful drain into a wait followed by
SIGKILL, with in-flight work lost --- exactly as if you had written no shutdown
handler at all.
\end{warning}

\subsection{Managed, or not}
\label{sec:platform}

LangGraph Platform provides the runtime this book has been building by hand:
threads, runs, cron, queueing, scaling, and the double-texting policies from
Chapter~\ref{ch:interrupts} as configuration.

The trade is the usual one, and worth stating plainly rather than resolving:
self-hosting means you own the per-thread lock, the queue, the drain and the
retention job --- all of which this book has shown you how to build, and none of
which is free to operate. Chapter~\ref{ch:ecosystem} weighs it properly.

% ============================================================
\section{Pinning, in an ecosystem that ships weekly}
\label{sec:pinning}

Chapter~\ref{ch:landscape} counted the cadence: dozens of releases a year across
the core packages, none of them breaking within 1.x. That combination argues for
a specific posture --- pin exactly in a committed lockfile, upgrade deliberately
on a schedule rather than reactively, and \textbf{gate the upgrade on the
evaluation set} (Ch.~\ref{ch:observability}).

That last part is what makes upgrading survivable. A minor version that changes
nothing in the API can still change behaviour, and the evaluation set is the only
thing that will tell you before your users do.

% ============================================================
\section{What to take away}
\label{sec:ch15-summary}

\begin{itemize}[leftmargin=1.4em]
\item Every effective control is enforced outside the model. Instructing it to
      behave is a preference.
\item Tools run with the user's permissions, checked in \code{wrap\_tool\_call}.
      Identity comes from context, never from state.
\item PII redaction is input-only by default.
\item \textbf{Checkpoints are personal data} --- retention, deletion, review.
      And a checkpoint is not an audit log.
\item The two large cost levers are context strategy and topology, both measured
      earlier in this book. The rest is enforcement.
\item Alert on the budget, not the invoice.
\item Syntax validation is not a security boundary. The container and the network
      policy are.
\item The grace period must exceed the longest run, and one swallowed
      cancellation defeats the whole drain.
\item Pin exactly; upgrade on a schedule; gate the upgrade on the evaluation set.
\end{itemize}

\begin{exercisebox}
\textbf{1.} Write a document containing an injection, put it where a tool will
fetch it, and watch a naive agent comply. Then add an allowlist and watch it
fail. Keep both outputs --- the first is what convinces a sceptical colleague.

\textbf{2.} Implement authorisation in \code{wrap\_tool\_call} and confirm the
tool body never executes for an unpermitted user.

\textbf{3.} Write the checkpoint retention job. Then answer, for your own
service, how you would satisfy a deletion request today.

\textbf{4.} Add a budget cap that halts a run mid-flight, and confirm the
checkpoint left behind is resumable.

\textbf{5.} Send SIGTERM to your service with a long run in flight. Time the
drain. Then swallow a \code{CancelledError} somewhere and time it again.
\end{exercisebox}

\begin{projectbox}
\textbf{Stage 12 --- security, cost, deployment.} Injection tests, authorisation
in the tool wrapper, PII redaction with a retention job, a budget cap, an audit
log distinct from the checkpoints, a Dockerfile, manifests and a SIGTERM handler.

The contract requires \textbf{both halves} of the injection test: a known
injection must \emph{succeed} against a naive configuration and \emph{fail}
against the hardened one. A test that only shows the hardened case passes proves
nothing about the hardening.
\end{projectbox}
